\chapter{2007年北京卷（文）}
\section{单选题}
\begin{problem}
已知 \(\cos \theta \cdot \tan \theta < 0\)，那么角 \(\theta\) 是（\quad）

\choices
    {第一或第二象限角}
    {第二或第三象限角}
    {第三或第四象限角}
    {第一或第四象限角}
\end{problem}

\begin{answer}
C.
\end{answer}

\begin{solution}
因为 \(\tan\theta=\dfrac{\sin\theta}{\cos\theta}\)，且题设中的式子有意义，所以 \(\cos\theta\neq0\). 于是
\(\cos\theta\cdot\tan\theta=\sin\theta<0\)，故 \(\theta\) 是第三或第四象限角，选 C.
\end{solution}

\begin{problem}
函数 \( f(x) = 3^{x} \)（\(0 < x \leqslant 2\)）的反函数的定义域为（\quad）

\choices
    {\((0, +\infty)\)}
    {\((1, 9]\)}
    {\((0,1)\)}
    {\([9, +\infty)\)}
\end{problem}

\begin{answer}
B.
\end{answer}

\begin{solution}
当 \(0<x\leqslant2\) 时，\(f(x)=3^x\) 的值域为 \(\left(3^0,3^2\right]=(1,9]\). 反函数的定义域等于原函数的值域，所以选 B.
\end{solution}

\begin{problem}
函数 \( f(x) = \sin 2x - \cos 2x \) 的最小正周期是（\quad）

\choices
    {\(\frac{\pi}{2}\)}
    {\(\pi\)}
    {\(2\pi\)}
    {\(4\pi\)}
\end{problem}

\begin{answer}
B.
\end{answer}

\begin{solution}
由
\(f(x)=\sin2x-\cos2x=\sqrt2\sin\left(2x-\dfrac{\pi}{4}\right)\)，可知其角频率为 \(2\)，所以最小正周期为 \(\dfrac{2\pi}{2}=\pi\)，选 B.
\end{solution}

\begin{problem}
椭圆 \(\frac{x^2}{a^2} + \frac{y^2}{b^2} = 1\)（\(a > b > 0\)）的焦点为 \(F_1\)，\(F_2\)，两条准线与 \(x\) 轴的交点分别为 \(M\)，\(N\). 若 \(|MN| \leqslant 2|F_1F_2|\)，则该椭圆离心率的取值范围是（\quad）

\choices
    {\(\left(0, \frac{1}{2}\right]\)}
    {\(\left(0, \frac{\sqrt{2}}{2}\right]\)}
    {\(\left[\frac{1}{2}, 1\right)\)}
    {\(\left[\frac{\sqrt{2}}{2}, 1\right)\)}
\end{problem}

\begin{answer}
D.
\end{answer}

\begin{solution}
设椭圆半焦距为 \(c\)，离心率为 \(e=\dfrac ca\). 两条准线为 \(x=\pm\dfrac{a}{e}\)，因此
\(|MN|=\dfrac{2a}{e}\)，而 \(|F_1F_2|=2c=2ae\). 由题设
\(\dfrac{2a}{e}\leqslant2\cdot2ae\)，得 \(e^2\geqslant\dfrac12\). 又 \(0<e<1\)，故 \(e\in\left[\dfrac{\sqrt2}{2},1\right)\)，选 D.
\end{solution}

\begin{problem}
某城市的汽车牌照号码由 \(2\) 个英文字母后接 \(4\) 个数字组成. 其中 \(4\) 个数字互不相同的牌照号码共有（\quad）

\choices
    {\(\left(\mathrm C_{26}^{1}\right)^{2}\mathrm A_{10}^{4}\) 个}
    {\(\mathrm A_{26}^{2}\mathrm A_{10}^{4}\) 个}
    {\(\left(\mathrm C_{26}^{1}\right)^{2}10^{4}\) 个}
    {\(\mathrm A_{26}^{2}10^{4}\) 个}
\end{problem}

\begin{answer}
A.
\end{answer}

\begin{solution}
两个英文字母各有 \(26\) 种选择，允许重复，故有 \(26^2=\left(\mathrm C_{26}^{1}\right)^2\) 种选法. 四个数字互不相同且有顺序，故有 \(\mathrm A_{10}^{4}\) 种选法. 所以牌照号码共有 \(\left(\mathrm C_{26}^{1}\right)^2\mathrm A_{10}^{4}\) 个，选 A.
\end{solution}

\begin{problem}
若不等式组 \( \begin{cases}
x-y+5&\geqslant0,\\
y&\geqslant a,\\
0&\leqslant x\leqslant2
\end{cases} \) 表示的平面区域是一个三角形，则 \(a\) 的取值范围是（\quad）

\choices
    {\(a < 5\)}
    {\(a \geqslant 7\)}
    {\(5 \leqslant a < 7\)}
    {\(a < 5\) 或 \(a \geqslant 7\)}
\end{problem}

\begin{answer}
C.
\end{answer}

\begin{solution}
不等式组等价于 \(0\leqslant x\leqslant2\)，且 \(a\leqslant y\leqslant x+5\). 当 \(a<5\) 时，直线 \(y=a\) 与两条边界 \(x=0\)，\(x=2\) 都相交，区域为四边形；当 \(a=5\) 时，左侧交点退化为 \((0,5)\)，区域为三角形；当 \(5<a<7\) 时，直线 \(y=a\) 与 \(y=x+5\) 交于 \((a-5,a)\)，并与 \(x=2\) 相交，区域仍为三角形；当 \(a\geqslant7\) 时，区域至多退化为一点或为空集. 因此 \(5\leqslant a<7\)，选 C.
\end{solution}

\begin{problem}
平面 \(\alpha\) \(\parallel\) 平面 \(\beta\) 的一个充分条件是（\quad）

\choices
    {存在一条直线 \(a\)，\(a \parallel \alpha\)，\(a \parallel \beta\)}
    {存在一条直线 \(a\)，\(a \subset \alpha\)，\(a \parallel \beta\)}
    {存在两条平行直线 \(a\)，\(b\)，\(a \subset \alpha\)，\(b \subset \beta\)，\(a \parallel \beta\)，\(b \parallel \alpha\)}
    {存在两条异面直线 \(a\)，\(b\)，\(a \subset \alpha\)，\(b \subset \beta\)，\(a \parallel \beta\)，\(b \parallel \alpha\)}
\end{problem}

\begin{answer}
D.
\end{answer}

\begin{solution}
若两平面不平行而相交，设交线为 \(l\). 取一条与 \(l\) 平行且不在两平面内的直线，可满足选项 A；在平面 \(\alpha\) 内取一条与 \(l\) 平行的直线，可满足选项 B；分别在 \(\alpha\)，\(\beta\) 内取两条与 \(l\) 平行的直线，可满足选项 C. 因而 A，B，C 都不是充分条件.

对选项 D，\(a\subset\alpha\) 且 \(a\parallel\beta\)，所以 \(a\) 不与 \(\beta\) 相交，只能在平面 \(\alpha\) 内与交线 \(l\) 平行；同理 \(b\parallel l\). 于是 \(a\parallel b\)，与题设两直线异面矛盾. 因此 \(\alpha\)，\(\beta\) 不可能相交，只能平行，选 D.
\end{solution}

\begin{problem}
对于函数

\begin{circlelist}
\item \( f(x) = |x + 2| \)；
\item \( f(x) = (x - 2)^2 \)；
\item \( f(x) = \cos (x - 2) \).
\end{circlelist}

判断如下两个命题的真假： 

命题甲： \( f(x+2) \) 是偶函数； 

命题乙： \( f(x) \) 在 \( (-\infty,2) \) 上是减函数，在 \( (2,+\infty) \) 上是增函数； 

能使命题甲，乙均为真的所有函数的序号是（\quad）

\choices
    {\circled{1}\circled{2}}
    {\circled{1}\circled{3}}
    {\circled{2}}
    {\circled{3}}
\end{problem}

\begin{answer}
C.
\end{answer}

\begin{solution}
对函数 \(\circled{1}\)，\(f(x+2)=|x+4|\) 不是偶函数，且 \(f(x)\) 的折点为 \(-2\)，不满足命题乙. 对函数 \(\circled{2}\)，\(f(x+2)=x^2\) 是偶函数，且 \(f(x)=(x-2)^2\) 在 \((-\infty,2)\) 上递减、在 \((2,+\infty)\) 上递增. 对函数 \(\circled{3}\)，\(f(x+2)=\cos x\) 是偶函数，但余弦函数在给出的两个无穷区间上都不保持题述的单调性. 因此只有函数 \(\circled{2}\) 满足两个命题，选 C.
\end{solution}

\section{填空题}
\begin{problem}
\(f'(x)\) 是 \(f(x) = \frac{1}{3} x^3 + 2x + 1\) 的导函数，则 \(f'(-1)\) 的值是\fillinblank{}.
\end{problem}

\begin{answer}
\(3\).
\end{answer}

\begin{solution}
由幂函数求导法则，\(f'(x)=x^2+2\). 因而 \(f'(-1)=(-1)^2+2=3\).
\end{solution}

\begin{problem}
若数列 \(\{a_{n}\}\) 的前 \(n\) 项和 \(S_{n} = n^{2} - 10n\)（\(n = 1\)，\(2\)，\(3\)，\(\dots\)），则此数列的通项公式为\fillinblank{}.
\end{problem}

\begin{answer}
\(2n-11\).
\end{answer}

\begin{solution}
当 \(n=1\) 时，\(a_1=S_1=1-10=-9=2\times1-11\). 当 \(n\geqslant2\) 时，
\(a_n=S_n-S_{n-1}=n^2-10n-\left((n-1)^2-10(n-1)\right)=2n-11\). 所以数列的通项公式为 \(a_n=2n-11\).
\end{solution}

\begin{problem}
已知向量 \(\bs{a} = (2,4)\)，\(\bs{b} = (1,1)\). 若向量 \(\bs{b} \perp (\bs{a} + \lambda \bs{b})\)，则实数 \(\lambda\) 的值是\fillinblank{}.
\end{problem}

\begin{answer}
\(-3\).
\end{answer}

\begin{solution}
由垂直条件，\(\bs b\cdot\left(\bs a+\lambda\bs b\right)=0\). 代入 \(\bs a=(2,4)\)，\(\bs b=(1,1)\)，得
\((1,1)\cdot(2+\lambda,4+\lambda)=6+2\lambda=0\)，所以 \(\lambda=-3\).
\end{solution}

\begin{problem}
在 \(\triangle ABC\) 中，若 \(\tan A = \frac{1}{3}\)，\(C = 150^{\circ}\)，\(BC = 1\)，则 \(AB =\fillinblank{}\).
\end{problem}

\begin{answer}
\(\frac{\sqrt{10}}{2}\).
\end{answer}

\begin{solution}
因为 \(A+B+C=180^{\circ}\) 且 \(C=150^{\circ}\)，所以 \(0<A<30^{\circ}\)，从而 \(A\) 为锐角. 由 \(\tan A=\dfrac13\)，得 \(\sin A=\dfrac{1}{\sqrt{10}}\). 根据正弦定理，
\(\dfrac{AB}{\sin C}=\dfrac{BC}{\sin A}\)，故
\(AB=\dfrac{\sin150^{\circ}}{\sin A}=\dfrac{1/2}{1/\sqrt{10}}=\dfrac{\sqrt{10}}2\).
\end{solution}

\begin{problem}
2002 年在北京召开的国际数学家大会，会标是以我国古代数学家赵爽的弦图为基础设计的. 弦图是由四个全等直角三角形与一个小正方形拼成的一个大正方形（如图）. 如果小正方形的面积为 \(1\)，大正方形的面积为 \(25\)，直角三角形中较小的锐角为 \(\theta\)，那么 \(\cos 2\theta\) 的值等于\fillinblank{}.

\texfigure[max width=0.28\linewidth]{img_repaint/2007/beijing_liberal/q13_fig1.tex}
\end{problem}

\begin{answer}
\(\frac{7}{25}\).
\end{answer}

\begin{solution}
设弦图中直角三角形的两条直角边分别为 \(u\)，\(v\)，斜边为大正方形的边长 \(5\). 小正方形边长为 \(1\)，故 \(|u-v|=1\)，且 \(u^2+v^2=25\). 于是
\(2uv=u^2+v^2-(u-v)^2=25-1=24\)，从而 \((u+v)^2=25+24=49\)，即 \(u+v=7\). 取 \(u>v\)，则较小锐角 \(\theta\) 的余弦为 \(\dfrac{u}{5}\)，所以
\(\cos2\theta=2\cos^2\theta-1=\dfrac{2u^2-25}{25}=\dfrac{u^2-v^2}{25}=\dfrac{(u-v)(u+v)}{25}=\dfrac7{25}\).
\end{solution}

\begin{problem}
已知函数 \(f(x)\)，\(g(x)\) 分别由下表给出

\begin{center}
\begin{tabular}{|c|c|c|c|}
\hline
\(x\) & \(1\) & \(2\) & \(3\) \\
\hline
\(f(x)\) & \(2\) & \(1\) & \(1\) \\
\hline
\end{tabular}
\qquad
\begin{tabular}{|c|c|c|c|}
\hline
\(x\) & \(1\) & \(2\) & \(3\) \\
\hline
\(g(x)\) & \(3\) & \(2\) & \(1\) \\
\hline
\end{tabular}
\end{center}

则 \( f[g(1)] \) 的值为\fillinblank{}；当 \( g[f(x)] = 2 \) 时，\(x =\)\fillinblank{}.
\end{problem}

\begin{answer}
\(1\)，\(1\).
\end{answer}

\begin{solution}
由表得 \(g(1)=3\)，所以 \(f[g(1)]=f(3)=1\). 又由表可知 \(g(t)=2\) 当且仅当 \(t=2\)，故 \(g[f(x)]=2\) 等价于 \(f(x)=2\). 查表得 \(f(x)=2\) 当且仅当 \(x=1\). 因此两空依次为 \(1\)，\(1\).
\end{solution}

\section{解答题}
\begin{problem}
记关于 \(x\) 的不等式 \(\frac{x - a}{x + 1} < 0\) 的解集为 \(P\)，不等式 \(|x - 1| \leqslant 1\) 的解集为 \(Q\).

\begin{enumerate}
\item 若 \(a = 3\)，求 \(P\).
\item 若 \(Q \subseteq P\)，求正数 \(a\) 的取值范围.
\end{enumerate}
\end{problem}

\begin{solution}
\begin{enumerate}
\item 当 \(a=3\) 时，分式的分子、分母分别在 \(x=3\)，\(x=-1\) 处变号. 由于分式小于零时分子、分母异号，且 \(x=-1\) 不在定义域内，所以 \(P=\left(-1,3\right)\).
\item 由 \(|x-1|\leqslant1\)，得 \(Q=[0,2]\). 因为 \(a>0\)，不等式 \(\dfrac{x-a}{x+1}<0\) 的解集为 \(P=\left(-1,a\right)\). 要使 \([0,2]\subset\left(-1,a\right)\)，只需且必须有 \(a>2\). 因此正数 \(a\) 的取值范围为 \((2,+\infty)\).
\end{enumerate}
\end{solution}

\begin{problem}
数列 \(\{a_{n}\}\) 中，\(a_{1} = 2\)，\(a_{n + 1} = a_{n} + cn\)（\(c\) 是常数，\(n = 1\)，\(2\)，\(3\)，\(\dots\)），且 \(a_{1}\)，\(a_{2}\)，\(a_{3}\) 成公比不为 1 的等比数列.

\begin{enumerate}
\item 求 \(c\) 的值.
\item 求 \(\{a_{n}\}\) 的通项公式.
\end{enumerate}
\end{problem}

\begin{solution}
\begin{enumerate}
\item 由递推关系，\(a_2=2+c\)，\(a_3=2+3c\). 因为 \(a_1\)，\(a_2\)，\(a_3\) 成等比数列，所以
\((2+c)^2=2(2+3c)\)，即 \(c(c-2)=0\). 当 \(c=0\) 时三项均为 \(2\)，公比为 \(1\)，不符合题意，故 \(c=2\).
\item 由递推关系，\(a_n=a_1+c\sum_{k=1}^{n-1}k=2+2\cdot\dfrac{n(n-1)}2=n^2-n+2\). 当 \(n=1\) 时该式也成立，所以 \(a_n=n^2-n+2\).
\end{enumerate}
\end{solution}

\begin{problem}
如图，在 \(\mathrm{Rt}\triangle AOB\) 中，\(\angle OAB = \frac{\pi}{6}\)，斜边 \(AB = 4\). \(\mathrm{Rt}\triangle AOC\) 可以通过 \(\mathrm{Rt}\triangle AOB\) 以直线 \(AO\) 为轴旋转得到，且二面角 \(B-AO-C\) 是直二面角. \(D\) 是 \(AB\) 的中点.

\begin{enumerate}
\item 求证：平面 \(COD\) \(\perp\) 平面 \(AOB\).
\item 求异面直线 \(AO\) 与 \(CD\) 所成角的大小.
\end{enumerate}

\texfigure[max width=0.45\linewidth]{img_repaint/2007/beijing_liberal/q17_fig1.tex}
\end{problem}

\begin{solution}
\begin{enumerate}
\item 因为 \(\triangle AOB\)，\(\triangle AOC\) 都是直角三角形，所以 \(AO\perp OB\)，\(AO\perp OC\). 二面角 \(B-AO-C\) 为直二面角，故 \(OB\perp OC\). 于是直线 \(OC\) 垂直于平面 \(AOB\) 内相交的两条直线 \(AO\)，\(OB\)，从而 \(OC\perp\text{平面 }AOB\). 又 \(OC\subset\text{平面 }COD\)，所以平面 \(COD\perp\) 平面 \(AOB\).
\item 在平面 \(AOB\) 内过 \(D\) 作 \(DE\perp OB\)，垂足为 \(E\)，并连接 \(CE\). 因为 \(D\) 是 \(AB\) 的中点，得 \(DE\parallel AO\)，\(DE=\dfrac{AO}{2}\)，\(OE=\dfrac{OB}{2}\). 在直角三角形 \(AOB\) 中，\(AO=AB\cos\dfrac\pi6=2\sqrt3\)，\(OB=AB\sin\dfrac\pi6=2\)，所以 \(DE=\sqrt3\)，\(OE=1\). 又 \(OC=OB=2\)，且 \(OC\perp OB\)，故 \(CE=\sqrt{OC^2+OE^2}=\sqrt5\). 由于 \(DE\parallel AO\)，\(\angle CDE\) 就是异面直线 \(AO\) 与 \(CD\) 所成的角；又 \(AO\) 同时垂直于 \(OC\) 和 \(OE\)，故 \(DE\perp\) 平面 \(OCE\)，从而 \(CE\perp DE\). 在直角三角形 \(CDE\) 中，\(\tan\angle CDE=\dfrac{CE}{DE}=\dfrac{\sqrt5}{\sqrt3}=\dfrac{\sqrt{15}}3\). 因此所求角为 \(\arctan\dfrac{\sqrt{15}}3\).
\end{enumerate}
\end{solution}

\begin{problem}
某条公共汽车线路沿线共有 \(11\) 个车站（包括起点站和终点站）. 在起点站开出的一辆公共汽车上有 \(6\) 位乘客，假设每位乘客在起点站之外的各个车站下车是等可能的，求：

\begin{enumerate}
\item 这 \(6\) 位乘客在互不相同的车站下车的概率.
\item 这 \(6\) 位乘客中恰有 \(3\) 人在终点站下车的概率.
\end{enumerate}
\end{problem}

\begin{solution}
\begin{enumerate}
\item 起点站之外共有 \(10\) 个车站，每位乘客有 \(10\) 种等可能的下车选择，基本事件总数为 \(10^6\). 六位乘客在互不相同的车站下车的有序选法数为 \(\mathrm A_{10}^{6}\)，所以所求概率为
\(\dfrac{\mathrm A_{10}^{6}}{10^6}=\dfrac{10\cdot9\cdot8\cdot7\cdot6\cdot5}{10^6}=\dfrac{189}{1250}\).
\item 先从六位乘客中选出在终点站下车的三人，有 \(\mathrm C_6^3\) 种选法；其余三人不能在终点站下车，各有 \(9\) 个车站可选. 所以所求概率为
\(\dfrac{\mathrm C_6^3\cdot9^3}{10^6}=\dfrac{20\cdot729}{10^6}=\dfrac{729}{50000}\).
\end{enumerate}
\end{solution}

\begin{problem}
如图，矩形 \(ABCD\) 的两条对角线相交于点 \(M(2,0)\)，\(AB\) 边所在直线的方程为 \(x - 3y - 6 = 0\)，点 \(T(-1,1)\) 在 \(AD\) 边所在直线上.

\begin{enumerate}
\item 求 \(AD\) 边所在直线的方程.
\item 求矩形 \(ABCD\) 外接圆的方程.
\item 若动圆 \(P\) 过点 \(N(-2,0)\)，且与矩形 \(ABCD\) 的外接圆外切，求动圆 \(P\) 的圆心的轨迹方程.
\end{enumerate}

\texfigure[max width=0.48\linewidth]{img_repaint/2007/beijing_liberal/q19_fig1.tex}
\end{problem}

\begin{solution}
\begin{enumerate}
\item 直线 \(AB\) 的斜率为 \(\dfrac13\)，所以 \(AD\perp AB\) 的斜率为 \(-3\). 又 \(T(-1,1)\) 在 \(AD\) 上，故 \(y-1=-3(x+1)\)，即 \(AD\) 边所在直线的方程为 \(3x+y+2=0\).
\item 联立 \(x-3y-6=0\) 与 \(3x+y+2=0\)，得 \(A=(0,-2)\). 矩形的两条对角线互相平分且相等，所以其交点 \(M(2,0)\) 是外接圆圆心. 因而半径 \(MA=\sqrt{(2-0)^2+(0+2)^2}=2\sqrt2\)，外接圆方程为 \(\left(x-2\right)^2+y^2=8\).
\item 设动圆圆心为 \(P(x,y)\)，则动圆半径为 \(PN\). 固定圆的圆心为 \(M(2,0)\)，半径为 \(2\sqrt2\)，外切条件给出 \(PM=PN+2\sqrt2\)，即 \(PM-PN=2\sqrt2\). 因此 \(P\) 到两个定点 \(M(2,0)\)，\(N(-2,0)\) 的距离之差为定值 \(2\sqrt2\)，轨迹是双曲线的左支. 其半焦距为 \(c=2\)，实半轴为 \(a=\sqrt2\)，虚半轴满足 \(b^2=c^2-a^2=2\). 所以轨迹方程为
\(\dfrac{x^2}{2}-\dfrac{y^2}{2}=1\)，并取左支 \(x\leqslant-\sqrt2\).
\end{enumerate}
\end{solution}

\begin{problem}
已知函数 \(y = kx\) 与 \(y = x^2 + 2\)（\(x \geqslant 0\)）的图象相交于 \(A(x_1, y_1)\)，\(B(x_2, y_2)\). \(l_1\)，\(l_2\) 分别是 \(y = x^2 + 2\)（\(x \geqslant 0\)）的图象在 \(A\)，\(B\) 两点的切线，\(M\)，\(N\) 分别是 \(l_1\)，\(l_2\) 与 \(x\) 轴的交点.

\begin{enumerate}
\item 求 \(k\) 的取值范围.
\item 设 \(t\) 为点 \(M\) 的横坐标，当 \(x_{1} < x_{2}\) 时，写出 \(t\) 以 \(x_{1}\) 为自变量的函数式，并求其定义域和值域.
\item 试比较 \(|OM|\) 与 \(|ON|\) 的大小，并说明理由（\(O\) 是坐标原点）.
\end{enumerate}
\end{problem}

\begin{solution}
\begin{enumerate}
\item 联立两图象的方程，得 \(x^2-kx+2=0\). 由于交点 \(A\)，\(B\) 对应两个不同的正根，需有 \(k^2-8>0\)，且两根之和 \(k>0\). 因而 \(k>2\sqrt2\).
\item 抛物线 \(y=x^2+2\) 在横坐标为 \(x_1\) 处的切线为 \(y=2x_1x-x_1^2+2\). 令 \(y=0\)，得点 \(M\) 的横坐标 \(t=\dfrac{x_1^2-2}{2x_1}=\dfrac{x_1}{2}-\dfrac1{x_1}\). 交点横坐标满足 \(x_1x_2=2\)，且 \(x_1<x_2\)，所以 \(0<x_1<\sqrt2\). 函数 \(t(x_1)\) 的导数为 \(\dfrac12+\dfrac1{x_1^2}>0\)，故定义域为 \((0,\sqrt2)\)，值域为 \((-\infty,0)\).
\item 同理，点 \(N\) 的横坐标为 \(\dfrac{x_2}{2}-\dfrac1{x_2}\). 由 \(x_1x_2=2\)，有 \(\dfrac{x_2}{2}-\dfrac1{x_2}=\dfrac1{x_1}-\dfrac{x_1}{2}=-\left(\dfrac{x_1}{2}-\dfrac1{x_1}\right)=-t\). 因此 \(|OM|=|t|=|ON|\).
\end{enumerate}
\end{solution}
